
\documentclass[12pt,a4paper]{report}

%================================================
% PACKAGES
%================================================
\usepackage[utf8]{inputenc}
\usepackage[T1]{fontenc}
\usepackage{lmodern}
\usepackage{geometry}
\usepackage{setspace}
\usepackage{graphicx}
\usepackage{amsmath,amssymb}
\usepackage{booktabs}
\usepackage{array}
\usepackage{longtable}
\usepackage{enumitem}
\usepackage{xcolor}
\usepackage{hyperref}
\usepackage{float}
\usepackage{listings}
\usepackage{fancyhdr}
\usepackage{titlesec}

%================================================
% PAGE SETTINGS
%================================================
\geometry{
    left=2.5cm,
    right=2.5cm,
    top=2.5cm,
    bottom=2.5cm
}

\onehalfspacing

\hypersetup{
    colorlinks=true,
    linkcolor=blue,
    urlcolor=blue,
    citecolor=blue,
    pdftitle={HHO-SMA Research Project Blueprint},
    pdfauthor={Research Project Documentation}
}

\pagestyle{fancy}
\fancyhf{}
\fancyhead[L]{HHO-SMA Research Project}
\fancyhead[R]{Project Blueprint}
\fancyfoot[C]{\thepage}

%================================================
% CODE SETTINGS
%================================================
\lstset{
    basicstyle=\ttfamily\small,
    breaklines=true,
    frame=single,
    numbers=left,
    numberstyle=\tiny,
    keywordstyle=\color{blue},
    commentstyle=\color{gray},
    stringstyle=\color{red},
    showstringspaces=false
}

%================================================
% DOCUMENT
%================================================
\begin{document}

%================================================
% TITLE PAGE
%================================================
\begin{titlepage}
    \centering

    \vspace*{2cm}

    {\Huge\bfseries HHO-SMA Research Project\par}

    \vspace{0.5cm}

    {\Large Complete Research and Development Blueprint\par}

    \vspace{1.5cm}

    \rule{\textwidth}{1pt}

    \vspace{0.5cm}

    {\Large
    Hybrid Harris Hawks Optimization and
    Slime Mould Algorithm
    \par}

    \vspace{0.5cm}

    {\large
    Application to Energy-Efficient Clustering
    and Routing in Wireless Sensor Networks
    \par}

    \vspace{0.5cm}

    \rule{\textwidth}{1pt}

    \vfill

    \begin{flushleft}
    \textbf{Project type:} PhD Research Project\\[0.3cm]
    \textbf{Implementation language:} Python\\[0.3cm]
    \textbf{Original source:} \texttt{hho\_sma\_optimize.m}\\[0.3cm]
    \textbf{Current implementation:}
    \texttt{hho\_sma\_original.py}
    \end{flushleft}

    \vfill

    {\large November 2026}

\end{titlepage}

\tableofcontents
\clearpage

%================================================
\chapter{Project Overview}
%================================================

\section{Project Identity}

The HHO-SMA project investigates the development of a
hybrid metaheuristic algorithm combining Harris Hawks
Optimization (HHO) and Slime Mould Algorithm (SMA).

The long-term objective is to apply this hybrid
optimization strategy to clustering and routing
optimization in Wireless Sensor Networks (WSNs).

The original research idea was developed in MATLAB
and is associated with the source file:

\begin{center}
\texttt{hho\_sma\_optimize.m}
\end{center}

The current development strategy is to reconstruct,
verify, and extend the algorithm entirely in Python.

\section{Research Direction}

The general research direction is:

\begin{quote}
Optimizing cluster formation and energy-aware routing
in Wireless Sensor Networks using a hybrid
bio-inspired optimization algorithm.
\end{quote}

The main research challenges include:

\begin{itemize}
    \item Energy-aware cluster-head selection.
    \item Balanced and unequal cluster formation.
    \item Reduction of energy consumption.
    \item Mitigation of energy holes and hot spots.
    \item Energy-aware multi-hop routing.
    \item Network lifetime maximization.
    \item Scalability under different network densities.
\end{itemize}

%================================================
\chapter{Research Objectives}
%================================================

\section{General Objective}

The main objective is to investigate whether combining
HHO and SMA can provide an effective optimization
strategy for energy-efficient WSN clustering and routing.

\section{Specific Objectives}

\begin{enumerate}
    \item Reconstruct the original HHO-SMA methodology.
    \item Implement the verified algorithm in Python.
    \item Validate the optimizer using mathematical
    benchmark functions.
    \item Compare HHO-SMA against standalone HHO and SMA.
    \item Develop an energy-aware WSN clustering model.
    \item Integrate multi-hop routing optimization.
    \item Evaluate network lifetime and energy efficiency.
\end{enumerate}

%================================================
\chapter{Hybrid HHO-SMA Methodology}
%================================================

\section{Harris Hawks Optimization}

HHO is a population-based metaheuristic inspired by
the cooperative hunting behavior of Harris hawks.

Its principal mechanisms include:

\begin{itemize}
    \item Exploration.
    \item Exploitation.
    \item Energy-dependent search transitions.
    \item Soft besiege.
    \item Hard besiege.
    \item Rapid dives and Levy-flight-based updates.
\end{itemize}

HHO is intended to provide the principal population
search framework.

\section{Slime Mould Algorithm}

SMA is inspired by the adaptive foraging behavior
of slime moulds.

Its optimization mechanisms include:

\begin{itemize}
    \item Fitness-dependent weighting.
    \item Adaptive search behavior.
    \item Oscillatory movement.
    \item Exploration and exploitation balancing.
    \item Refinement of promising solutions.
\end{itemize}

\section{Hybridization Concept}

The intended research concept is to integrate
SMA-inspired refinement into the HHO optimization
process.

The hybrid should not be assumed to consist of two
independent populations unless this is confirmed by
the original source.

The conceptual structure is:

\begin{equation}
\text{Initial Population}
\longrightarrow
\text{HHO Search}
\longrightarrow
\text{SMA Refinement}
\longrightarrow
\text{Fitness Evaluation}
\end{equation}

The exact update equations, transition conditions,
and integration mechanisms must be recovered from
the original MATLAB source before the method is
claimed to reproduce the original algorithm.

\section{Research Concepts to Preserve}

Potential research extensions include:

\begin{itemize}
    \item Energy-aware cluster-head selection.
    \item Refinement around energy-aware centroids.
    \item Unequal clustering.
    \item Energy-aware multi-hop routing.
    \item Adaptive optimization parameters.
\end{itemize}

These concepts are research directions and should
not automatically be considered part of the original
MATLAB implementation.

%================================================
\chapter{Python Development Environment}
%================================================

\section{Project Location}

The project is maintained under:

\begin{verbatim}
G:\My Drive\DOCTORAT\HHO_SMA_Project
\end{verbatim}

\section{Proposed Project Structure}

\begin{verbatim}
HHO_SMA_Project/
|
|-- algorithms/
|   |-- hho_sma_original.py
|
|-- benchmarks/
|   |-- sphere.py
|   |-- rastrigin.py
|   |-- ackley.py
|   |-- rosenbrock.py
|
|-- experiments/
|   |-- benchmark_experiments.py
|   |-- statistical_analysis.py
|   |-- wsn_experiments.py
|
|-- visualization/
|   |-- convergence.py
|   |-- network_topology.py
|   |-- energy_analysis.py
|
|-- results/
|   |-- benchmark_results/
|   |-- convergence/
|   |-- wsn_results/
|
|-- main.py
|-- requirements.txt
\end{verbatim}

This is the intended architecture. Not all modules
have been implemented yet.

\section{Software Requirements}

The implementation uses Python.

The principal dependencies are:

\begin{lstlisting}
numpy
matplotlib
scipy
pandas
\end{lstlisting}

Installation command:

\begin{lstlisting}
python -m pip install numpy matplotlib scipy pandas
\end{lstlisting}

\section{Execution}

From the project root:

\begin{lstlisting}
python -m algorithms.hho_sma_original
\end{lstlisting}

%================================================
\chapter{Current Experimental Status}
%================================================

\section{Sphere Benchmark}

The first experiment uses the Sphere function:

\begin{equation}
f(\mathbf{x}) =
\sum_{i=1}^{D}x_i^2
\end{equation}

where:

\begin{equation}
D=30
\end{equation}

The theoretical global optimum is:

\begin{equation}
f(\mathbf{x}^{*})=0
\end{equation}

with:

\begin{equation}
\mathbf{x}^{*}=(0,0,\ldots,0)
\end{equation}

\section{Experimental Parameters}

\begin{table}[H]
\centering
\caption{Current Sphere benchmark configuration}
\begin{tabular}{ll}
\toprule
\textbf{Parameter} & \textbf{Value}\\
\midrule
Population size & 30\\
Maximum iterations & 500\\
Dimension & 30\\
Lower bound & -100\\
Upper bound & 100\\
Benchmark function & Sphere\\
Implementation & Python\\
\bottomrule
\end{tabular}
\end{table}

\section{Recorded Results}

The current convergence experiment produced:

\begin{table}[H]
\centering
\caption{Best fitness at selected iterations}
\begin{tabular}{cc}
\toprule
\textbf{Iteration} & \textbf{Best Fitness}\\
\midrule
50  & $3.8878\times10^{-10}$\\
100 & $4.4057\times10^{-35}$\\
200 & $3.4430\times10^{-82}$\\
300 & $3.0461\times10^{-156}$\\
414 & 0\\
500 & 0\\
\bottomrule
\end{tabular}
\end{table}

Additional observations:

\begin{itemize}
    \item First recorded zero: iteration 414.
    \item Nonzero history entries: 413.
    \item Zero history entries: 87.
    \item Total iterations: 500.
\end{itemize}

The recorded convergence history is consistent with
progressive improvement followed by retention of the
best solution.

However, the final position and its Euclidean norm
must still be inspected to verify the solution.

\section{Output Files}

The current experiment saves:

\begin{verbatim}
results/sphere_convergence.csv
results/sphere_convergence.png
\end{verbatim}

%================================================
\chapter{Algorithm Verification}
%================================================

\section{Original MATLAB Source Audit}

The original MATLAB source file is:

\begin{center}
\texttt{hho\_sma\_optimize.m}
\end{center}

The next development stage must compare the MATLAB
implementation with the current Python implementation.

The audit must examine:

\begin{enumerate}
    \item Population initialization.
    \item Fitness evaluation.
    \item SMA weighting equations.
    \item HHO energy calculation.
    \item Exploration conditions.
    \item Exploitation conditions.
    \item Soft and hard besiege equations.
    \item Levy-flight updates.
    \item Greedy selection.
    \item Boundary handling.
    \item Termination criteria.
\end{enumerate}

Each equation must be classified as:

\begin{itemize}
    \item Directly derived from the original source.
    \item A Python implementation correction.
    \item A proposed methodological improvement.
\end{itemize}

\section{Final Position Verification}

The optimizer must return or expose the final
best position.

Its Euclidean norm should be calculated as:

\begin{equation}
\|\mathbf{x}^{*}\|_2 =
\sqrt{\sum_{i=1}^{D}(x_i^{*})^2}
\end{equation}

For the Sphere function, a zero objective value
should correspond to a zero position, subject to
floating-point precision.

%================================================
\chapter{Benchmark Validation}
%================================================

After verifying the Sphere experiment, the algorithm
should be evaluated on additional benchmark functions.

\begin{table}[H]
\centering
\caption{Planned benchmark functions}
\begin{tabular}{lll}
\toprule
\textbf{Function} & \textbf{Type} & \textbf{Purpose}\\
\midrule
Sphere & Unimodal & Basic convergence\\
Rosenbrock & Unimodal & Narrow valley search\\
Rastrigin & Multimodal & Local optima avoidance\\
Ackley & Multimodal & Exploration capability\\
Griewank & Multimodal & Complex search landscape\\
\bottomrule
\end{tabular}
\end{table}

\section{Independent Runs}

At least 30 independent runs should be performed
for the initial comparative study.

The following statistics should be reported:

\begin{itemize}
    \item Best fitness.
    \item Worst fitness.
    \item Mean fitness.
    \item Median fitness.
    \item Standard deviation.
    \item Success rate.
    \item Execution time.
\end{itemize}

\section{Baseline Algorithms}

The hybrid algorithm should be compared with:

\begin{itemize}
    \item Harris Hawks Optimization (HHO).
    \item Slime Mould Algorithm (SMA).
    \item Hybrid HHO-SMA.
\end{itemize}

All algorithms should use comparable experimental
budgets and identical benchmark conditions.

Statistical tests may include:

\begin{itemize}
    \item Wilcoxon signed-rank test.
    \item Friedman test.
    \item Appropriate post-hoc comparisons.
\end{itemize}

%================================================
\chapter{WSN Application}
%================================================

\section{Network Initialization}

The proposed WSN simulator should support:

\begin{itemize}
    \item Two-dimensional sensor deployment.
    \item Base station positioning.
    \item Initial sensor energy.
    \item Communication and aggregation parameters.
    \item Radio energy consumption.
    \item Packet transmission.
\end{itemize}

\section{Solution Representation}

A candidate solution may represent:

\begin{itemize}
    \item Cluster-head selection.
    \item Cluster membership.
    \item Unequal cluster sizes or radii.
    \item Routing decisions.
\end{itemize}

A suitable encoding and decoding mechanism must
be developed to transform optimizer solutions into
valid WSN configurations.

\section{Proposed Fitness Function}

A conceptual multi-objective-inspired scalar
fitness function is:

\begin{equation}
\begin{split}
F(X)={}&w_1E_{\mathrm{consumption}}
+w_2D_{\mathrm{intra}}\\
&+w_3D_{\mathrm{BS}}
+w_4B_{\mathrm{imbalance}}
+w_5P_{\mathrm{penalty}}
\end{split}
\end{equation}

where:

\begin{itemize}
    \item $E_{\mathrm{consumption}}$: energy cost.
    \item $D_{\mathrm{intra}}$: intra-cluster distance cost.
    \item $D_{\mathrm{BS}}$: routing cost to the base station.
    \item $B_{\mathrm{imbalance}}$: cluster-load imbalance.
    \item $P_{\mathrm{penalty}}$: feasibility penalty.
\end{itemize}

The weights and normalization must be experimentally
justified. This is a proposed design, not a finalized
fitness function.

\section{Network Lifetime Metrics}

The principal network lifetime indicators are:

\begin{description}
    \item[FND:] Round in which the first sensor node dies.
    \item[HND:] Round in which half of the sensor nodes die.
    \item[LND:] Round in which the last sensor node dies.
\end{description}

Additional metrics should include:

\begin{itemize}
    \item Number of alive nodes per round.
    \item Residual energy.
    \item Total energy consumption.
    \item Packet delivery ratio.
    \item Throughput.
    \item End-to-end delay.
    \item Control overhead.
\end{itemize}

\section{Baseline WSN Protocols}

Potential comparison protocols include:

\begin{itemize}
    \item LEACH.
    \item PEGASIS.
    \item Other verified energy-aware clustering methods.
\end{itemize}

All comparison protocols must use documented
and comparable simulation conditions.

%================================================
\chapter{Additional Research Extensions}
%================================================

The following ideas may be investigated after
the core hybrid algorithm has been validated:

\begin{enumerate}
    \item Unequal clustering for energy-hole mitigation.
    \item Energy-aware multi-hop routing.
    \item Refinement around energy-aware centroids.
    \item Stochastic duty cycling.
    \item Adaptive fitness weights.
    \item Residual-energy-based adaptation.
\end{enumerate}

These extensions should be evaluated individually
before combining them into a complete framework.

%================================================
\chapter{Development Roadmap}
%================================================

\begin{longtable}{p{0.12\textwidth}p{0.32\textwidth}p{0.48\textwidth}}
\toprule
\textbf{Phase} & \textbf{Activity} & \textbf{Expected Outcome}\\
\midrule
A & MATLAB/Python audit & Verified mathematical methodology\\
B & Sphere verification & Correct implementation and solution\\
C & Independent runs & Reproducibility and statistical results\\
D & Benchmark suite & General optimization validation\\
E & Baseline comparison & HHO, SMA and HHO-SMA evaluation\\
F & WSN simulation & Functional network model\\
G & Clustering optimization & Energy-aware cluster formation\\
H & Routing integration & Multi-hop energy-aware routing\\
I & Network experiments & FND, HND, LND and energy results\\
J & Research writing & Scientifically defensible publication\\
\bottomrule
\end{longtable}

%================================================
\chapter{Scientific and Development Rules}
%================================================

The following principles must be respected throughout
the project:

\begin{enumerate}
    \item Do not invent or misrepresent algorithm equations.
    \item Do not claim source fidelity without auditing
    the original MATLAB implementation.
    \item Do not silently modify the optimization method.
    \item Separate original methodology from improvements.
    \item Use reproducible random seeds.
    \item Preserve raw experimental results.
    \item Compare algorithms under fair conditions.
    \item Use real and verifiable scientific references.
    \item Do not claim superiority based on one experiment.
    \item Prioritize scientific validity over impressive results.
\end{enumerate}

%================================================
\chapter{Immediate Action Checklist}
%================================================

\begin{table}[H]
\centering
\caption{Current project status}
\begin{tabular}{p{0.72\textwidth}c}
\toprule
\textbf{Task} & \textbf{Status}\\
\midrule
Python project created & Done\\
Initial HHO-SMA prototype implemented & Done\\
Sphere benchmark executed & Done\\
CSV convergence history generated & Done\\
Progressive convergence observed & Done\\
Final position verification & Pending\\
Original MATLAB source audit & Pending\\
Python equation verification & Pending\\
Thirty independent runs & Pending\\
HHO and SMA comparison & Pending\\
Additional benchmark functions & Pending\\
WSN simulator & Pending\\
Clustering and routing integration & Pending\\
Network lifetime evaluation & Pending\\
\bottomrule
\end{tabular}
\end{table}

%================================================
\chapter{Instructions for the Next AI Agent}
%================================================

The next AI agent should continue this project
without assuming that the current Python implementation
faithfully reproduces the original MATLAB methodology.

The first task is to inspect:

\begin{enumerate}
    \item \texttt{hho\_sma\_optimize.m}
    \item \texttt{algorithms/hho\_sma\_original.py}
\end{enumerate}

It must reconstruct the original mathematical
methodology, compare the two implementations,
identify differences, and report the findings.

No major modifications should be made before
this verification.

The implementation must remain Python-based.
Complete code files and Windows PowerShell
instructions should be provided whenever code
changes are required.

The ultimate research objective is a scientifically
validated HHO-SMA optimizer for energy-efficient
clustering and routing in Wireless Sensor Networks.

%================================================
% END
%================================================

\end{document}
```

